\chapter{Bayes' Theorem and Reasoning Under Uncertainty}

There is a single, short equation in this chapter. You can write it on the back of a
business card, and it has been written on more than a few. Yet that one line is the
engine behind spam filters, medical diagnosis, courtroom reasoning, the search for lost
submarines, the calibration of climate models, and most of modern machine learning. It
is the rule that tells you, precisely and honestly, \emph{how a rational mind should
change its beliefs when new evidence arrives}.

That last sentence is worth pausing on, because it is the whole chapter in disguise.
Probability up to this point has mostly answered the question ``how likely is this
outcome?'' Bayes' theorem answers a different and far more human question: ``I observed
\emph{something}; what should I now believe about the hidden cause behind it?'' A doctor
sees a positive test and asks about the unseen disease. A detective sees a fingerprint
and asks about the unseen suspect. A phone's autocorrect sees the keystrokes
\texttt{teh} and asks about the unseen word you meant. In every case we are running a
chain of cause and effect \emph{backwards}: from the visible effect to the invisible
cause. Bayes' theorem is the mathematics of running that arrow backwards correctly.

We will derive it in one line, then spend the rest of the chapter learning to
\emph{feel} it. We will name its four parts --- prior, likelihood, evidence, posterior
--- and watch what each one does. We will meet the base-rate fallacy, perhaps the most
expensive single mistake in all of applied probability, and see exactly how an honest
test can produce mostly-wrong-looking results. We will finally put the Monty Hall
problem to rest with a method that leaves no room for argument. We will learn to update
not once but repeatedly, letting evidence accumulate. And we will end by building, by
hand, the algorithm that for two decades kept the spam out of your inbox.

\chapterline

\section{Running the Arrow Backwards}
\label{ch05:backwards}

Let us start with a story, because Bayes' theorem is at heart a story about asymmetry.

\begin{intuitionbox}
Imagine two facts about a far-off town. First: \emph{if it is raining, the streets are
wet} --- this happens almost always, say with probability $0.95$. Second: you are told,
by a friend who glanced out a window, only that \emph{the streets are wet}. How likely
is it that it is raining?

Your instinct might be to reach for that $0.95$. But notice the question has flipped.
You were given the chance of \emph{wet streets given rain}, and you were asked for the
chance of \emph{rain given wet streets}. These are not the same number, and confusing
them is the original sin of probabilistic reasoning. Streets can be wet for many reasons
--- a street cleaner, a burst pipe, melting snow, a parade with water balloons. To
answer the question you actually need a third fact you have not been given: \emph{how
often does it rain in this town at all?} In a rainforest the answer is ``almost surely
raining.'' In a desert, ``almost surely not, despite the wet streets.'' Same evidence,
opposite conclusions. The missing ingredient is what we will call the \emph{prior}, and
half of all reasoning errors come from forgetting it.
\end{intuitionbox}

The asymmetry in that story --- $\P(\text{wet} \given \text{rain})$ versus
$\P(\text{rain} \given \text{wet})$ --- is exactly the asymmetry Bayes' theorem
resolves. Recall from the previous chapter the definition of conditional probability:
for events $A$ and $B$ with $\P(B) > 0$,
\[
  \P(A \given B) \;=\; \frac{\P(A \cap B)}{\P(B)}.
\]
Read it as: ``of all the ways $B$ can happen, what fraction also have $A$?'' Now here is
the trick. The intersection $A \cap B$ does not care which letter we write first;
$\P(A \cap B) = \P(B \cap A)$. So we can expand that same joint probability two different
ways using the multiplication rule:
\[
  \P(A \cap B) \;=\; \P(A \given B)\,\P(B)
  \qquad\text{and}\qquad
  \P(A \cap B) \;=\; \P(B \given A)\,\P(A).
\]
Both equal the same thing, so they equal each other:
\[
  \P(A \given B)\,\P(B) \;=\; \P(B \given A)\,\P(A).
\]
Divide both sides by $\P(B)$ and you have arrived.

\begin{formalbox}
\textbf{Bayes' Theorem (basic form).} For events $A$ and $B$ with $\P(B) > 0$,
\begin{equation}
  \P(A \given B) \;=\; \frac{\P(B \given A)\,\P(A)}{\P(B)}.
  \label{ch05:bayes-basic}
\end{equation}
\end{formalbox}

That is the entire derivation. One symmetry of the intersection, one division. It is
almost embarrassing how little machinery it takes --- and yet the consequences fill this
chapter and several careers. The reason it feels deep even though it is shallow to prove
is that it tells you something your intuition resists: that to flip a conditional you
\emph{must} bring in the unconditional probabilities $\P(A)$ and $\P(B)$, and people
chronically leave them out.

\begin{whybox}
Why is $\P(A)$ in there at all? Because $\P(B\given A)$ only tells you how strongly the
cause $A$ \emph{predicts} the effect $B$ --- it says nothing about how common $A$ is to
begin with. A cause can be a perfect predictor of its effect and still be a terrible
explanation for it, simply because the cause is rare. ``Being struck by lightning''
predicts ``hospital visit'' very strongly, but almost no hospital visit is from
lightning, because lightning strikes are rare. The factor $\P(A)$ is the brake that
keeps a powerful-but-rare explanation from running away with your belief. Equation
\eqref{ch05:bayes-basic} is the precise bookkeeping of how predictive power
($\P(B\given A)$) and prior plausibility ($\P(A)$) trade off.
\end{whybox}

\subsection{The four names}

Equation \eqref{ch05:bayes-basic} earns its fame once we attach a story to each piece.
Let $A$ be a \emph{hypothesis} $H$ (the hidden cause: it is raining, the patient is
sick, the email is spam) and let $B$ be the \emph{evidence} $E$ we actually observe (wet
streets, a positive test, the word ``Viagra''). Rewriting,
\begin{equation}
  \underbrace{\P(H \given E)}_{\text{posterior}}
  \;=\;
  \frac{\overbrace{\P(E \given H)}^{\text{likelihood}}\;\overbrace{\P(H)}^{\text{prior}}}
       {\underbrace{\P(E)}_{\text{evidence}}}.
  \label{ch05:bayes-named}
\end{equation}

\begin{itemize}
\item The \textbf{prior} $\P(H)$ is what you believed about the hypothesis \emph{before}
seeing this piece of evidence. The rainfall rate of the town. Your starting credence.
\item The \textbf{likelihood} $\P(E \given H)$ is how well the hypothesis predicts the
evidence: \emph{if} $H$ were true, how probable would $E$ be? Note carefully: the
likelihood is a function of $H$, asking ``which hypothesis makes the data look least
surprising?'' It is \emph{not} a probability of $H$.
\item The \textbf{evidence} (or \emph{marginal likelihood}) $\P(E)$ is the total
probability of seeing $E$ at all, under every hypothesis weighted by its prior. It is
the great normalizer, the thing that makes the posterior add up to one.
\item The \textbf{posterior} $\P(H \given E)$ is your updated belief after folding in the
evidence. It is the answer. It is what you came for.
\end{itemize}

\begin{happeningbox}
Here is the one sentence to remember when the symbols blur:
\[
  \text{posterior} \;\propto\; \text{likelihood} \times \text{prior}.
\]
The denominator $\P(E)$ does not depend on the particular hypothesis $H$ --- it is the
same number whether you are asking about rain or street cleaners --- so it cannot change
\emph{which} hypothesis wins. It only rescales all the posteriors so they sum to one.
What actually does the work, what actually moves your belief, is the product
$\P(E\given H)\,\P(H)$. New belief is old belief reweighted by how well each hypothesis
explained what you just saw. Everything else is normalization. If you internalize only
this proportionality, you understand Bayes' theorem.
\end{happeningbox}

\subsection{Computing the evidence: the law of total probability}

We will constantly need $\P(E)$, and almost never is it handed to us directly. Instead
we compute it by splitting the world into the hypothesis and its complement (or into a
full set of competing hypotheses) and adding up the ways $E$ can occur. If $H$ and its
negation $H^c$ together cover all possibilities, then
\begin{equation}
  \P(E) \;=\; \P(E \given H)\,\P(H) \;+\; \P(E \given H^c)\,\P(H^c).
  \label{ch05:total-prob}
\end{equation}
This is the \emph{law of total probability}: the evidence can arrive either through the
hypothesis being true or through it being false, and we tally both routes. More
generally, if $H_1, \dots, H_k$ are mutually exclusive and exhaustive (a \emph{partition}
of the sample space), then $\P(E) = \sum_{i=1}^{k} \P(E \given H_i)\,\P(H_i)$, and Bayes'
theorem becomes
\begin{equation}
  \P(H_j \given E) \;=\; \frac{\P(E \given H_j)\,\P(H_j)}
                              {\sum_{i=1}^{k} \P(E \given H_i)\,\P(H_i)}.
  \label{ch05:bayes-partition}
\end{equation}
The denominator is just the numerator summed over all hypotheses --- which is why the
posteriors are guaranteed to add to one. Let us put the rain story to rest with real
numbers.

\begin{example}[Wet streets, finally answered]
\label{ch05:ex-rain}
Suppose in our town it rains on a given morning with probability $\P(R) = 0.2$. If it
rains, the streets are wet with probability $\P(W \given R) = 0.95$. If it does not rain,
the streets are still wet --- cleaners, sprinklers, the occasional spill --- with
probability $\P(W \given R^c) = 0.1$. Your friend reports wet streets. What is
$\P(R \given W)$?

First the evidence, by \eqref{ch05:total-prob}:
\[
  \P(W) = (0.95)(0.2) + (0.1)(0.8) = 0.19 + 0.08 = 0.27.
\]
Then Bayes:
\[
  \P(R \given W) = \frac{(0.95)(0.2)}{0.27} = \frac{0.19}{0.27} \approx 0.70.
\]
So wet streets push you from a $20\%$ belief in rain up to a $70\%$ belief --- a big
update, but nowhere near the $95\%$ that the careless reader grabbed from the likelihood.
The other $30\%$ of wet-street mornings are dry-but-wet for mundane reasons. Change the
town's rainfall rate to $\P(R)=0.02$ (a desert) and rerun it: $\P(W) = 0.019 + 0.098 =
0.117$ and $\P(R\given W) = 0.019/0.117 \approx 0.16$. Same evidence, same likelihoods,
but now wet streets leave you \emph{still doubting} the rain. The prior is doing exactly
the work the intuition box promised.
\end{example}

\begin{historybox}
The theorem is named for the Reverend \textbf{Thomas Bayes} (c.\ 1701--1761), a
Presbyterian minister and amateur mathematician in Tunbridge Wells, England. He never
published it. The result appeared posthumously in 1763 when his friend Richard Price
read Bayes' notes to the Royal Society in ``An Essay towards solving a Problem in the
Doctrine of Chances.'' Bayes' own problem was a charming one about a billiard ball rolled
on a table: having seen where later balls land relative to the first, what can you infer
about the first ball's hidden position? That is inference about a hidden cause from
visible effects --- exactly our theme.

But the person who truly built the cathedral was \textbf{Pierre-Simon Laplace}, who
around 1774, apparently unaware of Bayes, stated the rule in full generality and spent
decades applying it to astronomy, geodesy, and demography --- he used it to estimate,
from birth records, that boys are very slightly more likely than girls, and to put error
bars on the masses of the planets. For a century the method was simply called ``inverse
probability.'' It fell out of fashion in the early twentieth century, attacked by
frequentists like R.\ A.\ Fisher who distrusted the subjectivity of the prior, and was
kept alive in the wilderness by a few --- Jeffreys, de Finetti, Savage. It came roaring
back at the end of the century when cheap computation finally made the integrals
tractable, and today ``Bayesian'' is one of the most common adjectives in all of applied
science. A two-page note by a country parson became a worldview.
\end{historybox}

\section{The Base-Rate Fallacy: When a Good Test Lies}
\label{ch05:base-rate}

We now meet the single most important application of Bayes' theorem for everyday life,
and the place where untrained intuition fails most spectacularly and most expensively.
The setting is medical testing, but the lesson reaches into security screening,
fraud detection, courtroom DNA evidence, and every alarm system ever built.

\begin{paradoxbox}
\textbf{The disease that the test cannot find.} A certain disease afflicts $1$ person in
$1000$. There is a test for it that is, by any normal standard, excellent: if you have
the disease, the test is positive $99\%$ of the time (only a $1\%$ false-negative rate);
if you do \emph{not} have the disease, the test is negative $95\%$ of the time (a $5\%$
false-positive rate). You take the test. It comes back \textbf{positive}. How worried
should you be --- what is the probability you actually have the disease?

Most people, including a famously large fraction of doctors when surveyed, answer
somewhere around $95\%$ or $99\%$. The test is $95$--$99\%$ accurate, after all. The
true answer is about \textbf{2\%}. You are almost certainly fine. Read that again,
because it is genuinely shocking: a $99\%$-sensitive, $95\%$-specific test, returning
positive, and you are still $98\%$ likely to be healthy.
\end{paradoxbox}

How can an excellent test be so misleading? Bayes' theorem will tell us exactly, but
first let us see it with our bare hands, by counting people. This ``natural frequencies''
trick is the single best way to defuse the paradox, and it is worth doing even when you
trust the algebra.

\begin{walkbox}
Imagine $100{,}000$ people, a small city, all getting tested.

\emph{Step 1 --- the truly sick.} The disease rate is $1$ in $1000$, so about
$100{,}000 / 1000 = 100$ people genuinely have it. Of these $100$, the test catches
$99\%$, so $99$ of them test positive (and $1$ slips through as a false negative).

\emph{Step 2 --- the healthy many.} The other $99{,}900$ people are disease-free. The
test wrongly flags $5\%$ of them, so $0.05 \times 99{,}900 \approx 4995$ healthy people
also test positive.

\emph{Step 3 --- who got the bad news?} Total positives:
$99 + 4995 = 5094$. Of these, only $99$ are actually sick. So the probability that a
positive result means real disease is
\[
  \frac{99}{5094} \approx 0.0194 \approx 1.9\%.
\]
There is the answer, with no formula in sight. The flood of false positives from the
enormous healthy population (nearly $5000$ of them) drowns out the trickle of true
positives from the tiny sick population (just $99$). The test is fine. The problem is
that the disease is \emph{rare}, and rarity is the prior.
\end{walkbox}

Now the same computation in the language of Bayes, so you see they agree. Let $D$ be
``has disease'' and $+$ be ``tests positive.'' We are given the prior $\P(D) = 0.001$,
the likelihood (sensitivity) $\P(+ \given D) = 0.99$, and the false-positive rate
$\P(+ \given D^c) = 0.05$. The evidence:
\[
  \P(+) = \P(+\given D)\P(D) + \P(+\given D^c)\P(D^c)
        = (0.99)(0.001) + (0.05)(0.999)
        = 0.00099 + 0.04995 = 0.05094.
\]
The posterior:
\[
  \P(D \given +) = \frac{\P(+\given D)\P(D)}{\P(+)}
                 = \frac{0.00099}{0.05094} \approx 0.0194.
\]
Identical to the counting, as it must be --- the natural-frequencies method \emph{is}
Bayes' theorem with the denominators cleared.

\begin{trapbox}
The fallacy has a precise shape: people silently swap $\P(D \given +)$ for
$\P(+ \given D)$. They hear ``the test is $99\%$ accurate'' (a statement about
$\P(+\given D)$, the likelihood) and report it as the chance of being sick (the
posterior $\P(D\given +)$). These differ by exactly the factor $\P(D)/\P(+)$ that Bayes
inserts --- and when the disease is rare, that factor is tiny. Whenever you catch
yourself reasoning ``the evidence is very reliable, therefore the conclusion is very
likely,'' stop and ask: \emph{how common was the conclusion to begin with?} A reliable
detector of a rare thing is mostly an alarm that cries wolf.
\end{trapbox}

\begin{aestheticbox}
This is not an academic curiosity; it is a matter of public policy and personal anguish.
It is the central argument against mass screening of the general population for rare
diseases: you generate a torrent of frightened, falsely-positive patients --- biopsies,
anxiety, follow-up procedures with their own risks --- to find a handful of true cases.
It is why doctors test \emph{symptomatic} or \emph{high-risk} patients: doing so raises
the prior $\P(D)$, which is the only lever that improves the posterior. The same
arithmetic governs airport security, mass surveillance for terrorism, and automated fraud
flags: when you hunt for needles, almost everything your detector lights up is hay, no
matter how good the detector. The remedy is never ``a better test'' alone --- it is a
higher prior, achieved by testing a population in which the thing is not so rare.
\end{aestheticbox}

\subsection{Watching the prior take over}

The whole drama lives in how the posterior depends on the prior, the base rate. Let us
hold the test fixed (sensitivity $0.99$, specificity $0.95$) and sweep the disease
prevalence from very rare to fairly common, plotting $\P(D \given +)$ as we go. The
curve in Figure~\ref{ch05:fig-baserate} tells the whole story at a glance.

\begin{figure}[H]
\centering
\begin{tikzpicture}
\begin{axis}[
    width=12cm, height=7.2cm,
    xlabel={Disease prevalence (prior) $\P(D)$},
    ylabel={$\P(D \given +)$ after a positive test},
    xmin=0, xmax=0.5, ymin=0, ymax=1,
    domain=0.0005:0.5, samples=160,
    grid=both, grid style={gray!18},
    legend pos=south east, legend cell align=left,
    tick label style={font=\small}, label style={font=\small},
    title={A fixed test (sens.\ $0.99$, spec.\ $0.95$): the prior runs the show}
]
\addplot[densityred, very thick] {(0.99*x)/(0.99*x + 0.05*(1-x))};
\addlegendentry{posterior $\P(D\given +)$}
\addplot[sourceblue, dashed, thick] {x};
\addlegendentry{prior $=$ posterior (no info line)}
\addplot[only marks, mark=*, mark size=2pt, structurepurple]
    coordinates {(0.001, 0.0194)};
\node[anchor=west, font=\footnotesize, structurepurple]
    at (axis cs:0.02,0.10) {$1$ in $1000 \Rightarrow 1.9\%$};
\end{axis}
\end{tikzpicture}
\caption{For a strong but imperfect test, the post-positive probability of disease climbs
steeply with the base rate. At a prevalence of $1$ in $1000$ (purple dot) a positive
result barely moves you off zero; only when the disease is already reasonably common does
a positive test become trustworthy. The dashed line marks where evidence would tell you
nothing; the gap above it is how much the test \emph{does} help. Same test, wildly
different verdicts --- the only thing that changed is the prior.}
\label{ch05:fig-baserate}
\end{figure}

The shape of that curve is the base-rate fallacy made visible. Down in the far left,
where real diseases live, the red posterior hugs the floor: a positive test is almost
worthless because false positives swamp it. The posterior only becomes large --- the test
only becomes trustworthy --- once the prior is already substantial. The test never
changes; the verdict does. This is why the very same positive mammogram means something
very different for a $40$-year-old with no family history than for a $65$-year-old with a
suspicious lump: their priors are not the same, so their posteriors cannot be either.

\section{Sequential Updating: Yesterday's Posterior Is Today's Prior}
\label{ch05:sequential}

So far we have updated once. But evidence rarely arrives all at once; it dribbles in. The
beauty of Bayes' theorem is that it handles a stream of evidence with no extra machinery
at all. You simply update, and then update again, treating the belief you just computed as
the starting belief for the next round. This is the rhythm of all learning:

\begin{center}
prior $\xrightarrow{\;E_1\;}$ posterior$_1$ $=$ prior$_2$
$\xrightarrow{\;E_2\;}$ posterior$_2$ $=$ prior$_3$
$\xrightarrow{\;E_3\;}$ $\cdots$
\end{center}

\begin{whybox}
Why is it legitimate to feed yesterday's posterior in as today's prior? Because that is
literally what conditioning on more data does. Apply Bayes twice with two pieces of
evidence $E_1, E_2$:
\[
  \P(H \given E_1, E_2) = \frac{\P(E_2 \given H, E_1)\,\P(H \given E_1)}{\P(E_2\given E_1)}.
\]
The term $\P(H \given E_1)$ sitting where the prior usually sits is precisely the
posterior from round one. The data simply accumulate inside the conditioning bar. If the
pieces of evidence are conditionally independent given $H$ --- meaning that once you know
the true hypothesis, $E_1$ tells you nothing further about $E_2$ --- then
$\P(E_2 \given H, E_1) = \P(E_2 \given H)$, and the whole stream collapses to multiplying
likelihoods:
\[
  \P(H \given E_1, \dots, E_n) \;\propto\; \P(H)\prod_{i=1}^{n}\P(E_i \given H).
\]
Each new datum just multiplies in another likelihood factor. Belief compounds like
interest.
\end{whybox}

A spectacular feature of this process is its \emph{commutativity and convergence}: the
order in which conditionally-independent evidence arrives does not change the final
belief, and as evidence piles up the posterior typically concentrates on the truth almost
regardless of where the prior started (as long as the prior didn't dogmatically assign
probability zero to the truth --- a hypothesis you give zero prior can never be revived,
which is a moral as much as a theorem). Let us watch it happen.

\begin{example}[Two tests are much better than one]
\label{ch05:ex-second-test}
Return to the patient from the base-rate paradox, who tested positive and now stands at
posterior $\P(D \given +) \approx 0.0194$. The doctor, sensibly, orders a \emph{second},
independent test of the same quality. It too comes back positive. Now what?

Yesterday's posterior becomes today's prior: $\P(D) = 0.0194$. Run Bayes again with the
same likelihoods:
\[
  \P(+) = (0.99)(0.0194) + (0.05)(0.9806) = 0.01921 + 0.04903 = 0.06824,
\]
\[
  \P(D \given +,+) = \frac{(0.99)(0.0194)}{0.06824} = \frac{0.01921}{0.06824}
  \approx 0.281.
\]
Two positives lift the patient from $1.9\%$ to about $28\%$. A \emph{third} positive,
updating again with prior $0.281$, gives
\[
  \P(D\given +,+,+) = \frac{(0.99)(0.281)}{(0.99)(0.281)+(0.05)(0.719)}
  = \frac{0.2782}{0.2782+0.03595} \approx 0.886.
\]
Now you should be worried. This is exactly how diagnosis actually works: no single test
is conclusive for a rare disease, but independent confirmations compound. Each positive
multiplies the odds by the same factor; pile up enough of them and even a tiny prior is
overwhelmed.
\end{example}

Figure~\ref{ch05:fig-sequential} draws that compounding. Each positive test multiplies the
\emph{odds} by the constant Bayes factor $19.8$, so the posterior probability follows the
unmistakable S-shaped march from near-zero toward near-certainty: slow while the rare prior
still dominates, explosive once enough evidence has accumulated, then flattening as it
saturates against the ceiling of $1$.

\begin{figure}[H]
\centering
\begin{tikzpicture}
\begin{axis}[
    width=12cm, height=7cm,
    xlabel={Number of independent positive tests $n$},
    ylabel={posterior $\P(D\given n\text{ positives})$},
    xmin=0, xmax=8, ymin=0, ymax=1,
    xtick={0,1,2,3,4,5,6,7,8},
    ytick={0,0.2,0.4,0.6,0.8,1.0},
    grid=both, grid style={gray!18},
    tick label style={font=\small}, label style={font=\small},
    title={Each positive multiplies the odds by $19.8$: a tiny prior, overwhelmed}
]
% prior odds 1/999, Bayes factor 19.8^n; prob = o/(1+o), o = 19.8^n/999
\addplot[densityred, very thick, mark=*, mark size=2pt]
  coordinates {
    (0, 0.001000)
    (1, 0.019435)
    (2, 0.281832)
    (3, 0.885977)
    (4, 0.993542)
    (5, 0.999672)
    (6, 0.999983)
    (7, 0.999999)
    (8, 1.000000)
  };
\addplot[sourceblue, dashed, thick, domain=0:8] {0.5};
\node[anchor=south west, font=\footnotesize, sourceblue] at (axis cs:0.1,0.5) {coin-flip line};
\node[anchor=west, font=\footnotesize, structurepurple] at (axis cs:3.05,0.78)
  {3 positives $\Rightarrow 88.6\%$};
\end{axis}
\end{tikzpicture}
\caption{Repeated positive results on independent copies of the same imperfect test,
starting from a prevalence of $1$ in $1000$. The first positive barely registers (the
base-rate fallacy in action); the second and third do the real work; by the fourth the
patient is past the coin-flip line and climbing toward certainty. Belief compounds like
interest, because in log-odds each positive simply \emph{adds} the same fixed weight
$\log 19.8$.}
\label{ch05:fig-sequential}
\end{figure}

\subsection{The odds form: Bayes' theorem with the friction removed}

Example~\ref{ch05:ex-second-test} hints at something cleaner. Recomputing $\P(E)$ every
round is tedious, and it obscures the pattern. There is a reformulation in terms of
\emph{odds} that makes sequential updating almost trivial, and that reveals what evidence
``weighs.'' Take Bayes' theorem for $H$ and for $H^c$ and divide one by the other; the
shared denominator $\P(E)$ cancels:
\begin{equation}
  \underbrace{\frac{\P(H\given E)}{\P(H^c\given E)}}_{\text{posterior odds}}
  \;=\;
  \underbrace{\frac{\P(E\given H)}{\P(E\given H^c)}}_{\text{likelihood ratio}}
  \;\times\;
  \underbrace{\frac{\P(H)}{\P(H^c)}}_{\text{prior odds}}.
  \label{ch05:odds-form}
\end{equation}
This is the most useful single equation for actually \emph{doing} Bayesian updates by
hand. Odds are just probabilities written as a ratio: a probability of $0.2$ is odds of
$0.2/0.8 = 1{:}4$. The \emph{likelihood ratio} (also called the Bayes factor) is the
weight of the evidence --- how many times more likely the data are under $H$ than under
$H^c$. And the rule could not be simpler:
\[
  \text{posterior odds} \;=\; \text{Bayes factor} \times \text{prior odds}.
\]
No normalizer to compute. To chain evidence, just keep multiplying by each Bayes factor.

\begin{example}[The same diagnosis, the easy way]
For our test, the Bayes factor of a positive result is
\[
  \frac{\P(+\given D)}{\P(+\given D^c)} = \frac{0.99}{0.05} = 19.8.
\]
Every positive multiplies the odds of disease by $19.8$. The prior odds at prevalence
$1$ in $1000$ are $0.001/0.999 \approx 1/999$. After one positive:
$\frac{1}{999}\times 19.8 = \frac{19.8}{999} \approx 0.0198$, which is odds, i.e.\
probability $0.0198/(1.0198) \approx 0.0194$ --- matching before. After three positives:
$\frac{1}{999}\times 19.8^3 = \frac{7762}{999}\approx 7.77$, i.e.\ probability
$7.77/8.77 \approx 0.886$. Three multiplications and we are done, no total-probability
sums anywhere. \emph{This} is why practitioners think in odds and log-odds.
\end{example}

\begin{connectionbox}
Take the logarithm of the odds form \eqref{ch05:odds-form} and the products become sums:
\[
  \log\frac{\P(H\given E_{1:n})}{\P(H^c\given E_{1:n})}
  = \log\frac{\P(H)}{\P(H^c)} + \sum_{i=1}^{n}\log\frac{\P(E_i\given H)}{\P(E_i\given H^c)}.
\]
Each piece of evidence contributes an additive ``weight'' $\log(\text{Bayes factor})$ to
a running total of \emph{log-odds}. It is no accident that this looks like a linear model.
The logistic-regression classifier of machine learning is exactly this equation: it
learns a weight for each feature and \emph{adds them up} to form a log-odds, then squashes
back to a probability with the logistic function $\sigma(t) = 1/(1+e^{-t})$, which is
precisely the inverse of the log-odds transform. Information theory will later measure
these weights in \emph{bits} and \emph{nats} --- the ``weight of evidence'' is literally
a quantity of information. The thread from Reverend Bayes to your linear classifier runs
straight through this one logarithm.
\end{connectionbox}

The bridge between the additive world of log-odds and the bounded world of probabilities
is the logistic function $\sigma$, drawn in Figure~\ref{ch05:fig-logistic}. Adding up the
weights of evidence is easy and unbounded; the squashing curve then folds that running
total back into a probability in $[0,1]$, gently near the middle and ever more compressed
out in the tails --- which is exactly why the sequential-test curve of
Figure~\ref{ch05:fig-sequential} flattens as it approaches certainty.

\begin{figure}[H]
\centering
\begin{tikzpicture}
\begin{axis}[
    width=12cm, height=6.4cm,
    xlabel={log-odds $t$ (total weight of evidence)},
    ylabel={probability $\sigma(t)=\tfrac{1}{1+e^{-t}}$},
    xmin=-6, xmax=6, ymin=0, ymax=1,
    domain=-6:6, samples=160,
    xtick={-6,-4,-2,0,2,4,6},
    ytick={0,0.25,0.5,0.75,1.0},
    grid=both, grid style={gray!18},
    tick label style={font=\small}, label style={font=\small},
    title={From log-odds to probability: the logistic squashing function}
]
\addplot[densityred, very thick] {1/(1+exp(-x))};
\addplot[sourceblue, dashed, thick, domain=-6:6] {0.5};
\addplot[only marks, mark=*, mark size=2pt, structurepurple] coordinates {(0,0.5)};
\node[anchor=south west, font=\footnotesize, structurepurple] at (axis cs:0.15,0.5)
  {$t=0 \Rightarrow$ even odds};
\end{axis}
\end{tikzpicture}
\caption{The logistic function maps the running total of log-odds back to a probability.
Zero log-odds (evidence perfectly balanced) maps to probability $\tfrac12$; each unit of
accumulated evidence-weight slides you along the curve, with diminishing effect once you
are already near $0$ or $1$. Sequential Bayesian updating walks rightward along this curve,
one Bayes factor at a time.}
\label{ch05:fig-logistic}
\end{figure}

\section{The Monty Hall Problem, Finally Resolved}
\label{ch05:monty}

No book on Bayesian reasoning can skip the problem that broke the internet of its day.
It is a perfect specimen: a question simple enough for a child, an answer that PhDs got
wrong in print, and a resolution that Bayes' theorem delivers without drama.

\begin{paradoxbox}
\textbf{Monty Hall.} On a game show there are three doors. Behind one is a car; behind the
other two, goats. You pick a door --- say Door $1$ --- but do not open it. The host, Monty,
who knows where the car is, opens one of the \emph{other} two doors to reveal a goat ---
say he opens Door $3$. He then offers you a choice: stick with Door $1$, or switch to the
remaining closed Door $2$. \emph{Should you switch?}

The overwhelming intuition is: it doesn't matter. Two doors remain, one has the car, so
it's $50/50$ --- switching is pointless. This intuition is \textbf{wrong}. Switching wins
the car with probability $2/3$; staying wins with only $1/3$. You should always switch,
and doing so \emph{doubles} your chance of the car. When the columnist Marilyn vos Savant
published this answer in 1990, she received thousands of angry letters, many from
mathematicians, insisting she was wrong. She was not.
\end{paradoxbox}

\begin{historybox}
The problem is named for \textbf{Monty Hall}, host of the American television show
\emph{Let's Make a Deal}, though Monty's real game was looser than the puzzle. Its
mathematical skeleton dates to Joseph Bertrand's box paradox (1889) and to a 1975
letter by statistician Steve Selvin. It exploded into public consciousness in September
1990 when \textbf{Marilyn vos Savant} answered it in her \emph{Parade} magazine column.
The ensuing furor --- roughly $10{,}000$ letters, a sizable fraction from people with
doctorates, nearly all of them wrong and many of them rude --- became a landmark in the
history of how badly trained intuition handles conditional probability. Even Paul
Erd\H{o}s, one of the great mathematicians of the century, reportedly refused to believe
the answer until shown a computer simulation. If it fooled Erd\H{o}s, you are in good
company.
\end{historybox}

\subsection{Why the intuition fails, in words}

Before any algebra, here is the idea that makes it click for most people. Your original
door had a $1/3$ chance of hiding the car, and \emph{nothing Monty does afterward can
change a probability you already locked in}. So Door $1$ stays at $1/3$ forever. But the
total must be $1$, and Door $3$ is now revealed to be $0$. The entire remaining $2/3$ has
nowhere to go but onto Door $2$. Monty, by deliberately opening a goat door he \emph{chose
knowing the layout}, has effectively concentrated the other two doors' combined $2/3$
chance onto the single door he left closed.

\begin{intuitionbox}
If three doors still feels slippery, scale it up. Imagine \textbf{$100$ doors}, one car.
You pick Door $1$ --- a $1\%$ shot. Now Monty, who knows where the car is, opens
$98$ of the other doors, every one a goat, leaving just your Door $1$ and one other door,
say Door $57$, closed. Do you really believe it's now $50/50$? Of course not. You almost
certainly picked a goat at the start ($99\%$ chance), and Monty has carefully cleared away
every other goat, leaving the car sitting behind Door $57$ in $99$ cases out of $100$.
Switching wins $99\%$ of the time. The three-door case is the same effect, merely too
small to feel: staying wins $1/3$, switching wins $2/3$. The host's knowledge is the
whole trick --- he never opens the car, so his choice of door leaks information.
\end{intuitionbox}

\subsection{The Bayesian computation, leaving no doubt}

Words convince most people; algebra convinces everyone. Let us nail it with Bayes'
theorem. Define the hypotheses --- the car's true location:
\[
  C_1 = \{\text{car behind Door 1}\},\quad
  C_2 = \{\text{car behind Door 2}\},\quad
  C_3 = \{\text{car behind Door 3}\}.
\]
By symmetry before anything happens, $\P(C_1)=\P(C_2)=\P(C_3)=1/3$; these are our priors.
You have picked Door $1$. The evidence we will condition on is $M_3 = \{\text{Monty opens
Door 3 to show a goat}\}$. We need the likelihoods $\P(M_3 \given C_i)$ --- and \emph{here}
is where the host's rules, which the careless solver ignores, enter the math.

\begin{itemize}
\item $\P(M_3 \given C_1)$: the car is behind your Door $1$, so both other doors hide
goats and Monty may open either Door $2$ or Door $3$. With no reason to prefer one, he
opens Door $3$ with probability $\tfrac{1}{2}$. So $\P(M_3\given C_1)=\tfrac12$.
\item $\P(M_3 \given C_2)$: the car is behind Door $2$. Monty cannot open your Door $1$,
and he cannot open Door $2$ (the car!), so he is \emph{forced} to open Door $3$. Thus
$\P(M_3\given C_2)=1$.
\item $\P(M_3 \given C_3)$: the car is behind Door $3$. Monty will never open the car's
door, so he certainly does not open Door $3$. Thus $\P(M_3\given C_3)=0$.
\end{itemize}

Now assemble the evidence by total probability:
\[
  \P(M_3) = \tfrac12\cdot\tfrac13 + 1\cdot\tfrac13 + 0\cdot\tfrac13
          = \tfrac16 + \tfrac13 = \tfrac12.
\]
(Sensible: by symmetry Monty opens Door $3$ half the time.) Then the two posteriors that
matter:
\[
  \P(C_1 \given M_3) = \frac{\P(M_3\given C_1)\P(C_1)}{\P(M_3)}
  = \frac{\tfrac12\cdot\tfrac13}{\tfrac12} = \frac{\tfrac16}{\tfrac12} = \frac13,
\]
\[
  \P(C_2 \given M_3) = \frac{\P(M_3\given C_2)\P(C_2)}{\P(M_3)}
  = \frac{1\cdot\tfrac13}{\tfrac12} = \frac{\tfrac13}{\tfrac12} = \frac23.
\]
Staying with Door $1$ wins with probability $1/3$; switching to Door $2$ wins with
probability $2/3$. The paradox is dissolved, exactly, with no appeal to intuition.

The whole computation is laid bare by the probability tree in Figure~\ref{ch05:fig-monty},
which simply enumerates the worlds and weights each by its probability. Conditioning on
``Monty opened Door $3$'' means keeping only the branches consistent with that observation
and renormalizing --- and when you do, the surviving weight on ``car behind Door $2$'' is
twice the weight on ``car behind Door $1$.''

\begin{figure}[H]
\centering
\begin{tikzpicture}[
    grow=right, level distance=34mm, sibling distance=15mm,
    edge from parent/.style={draw, -{Latex[length=2mm]}, gray!70},
    every node/.style={font=\small, align=center}
]
\node[structurepurple] {start}
  child {node[sourceblue] {$C_3$ ($\tfrac13$)}
    child {node[gray!55] {opens 2 \\ (forced)} edge from parent node[below, font=\footnotesize]{$1$}}
    edge from parent node[above, font=\footnotesize]{$\tfrac13$}
  }
  child {node[sourceblue] {$C_2$ ($\tfrac13$)}
    child {node[densityred] {\textbf{opens 3} \\ $\Rightarrow \tfrac13$} edge from parent node[below, font=\footnotesize]{$1$}}
    edge from parent node[above, font=\footnotesize]{$\tfrac13$}
  }
  child {node[sourceblue] {$C_1$ ($\tfrac13$)}
    child {node[gray!55] {opens 2} edge from parent node[below, font=\footnotesize]{$\tfrac12$}}
    child {node[densityred] {\textbf{opens 3} \\ $\Rightarrow \tfrac16$} edge from parent node[above, font=\footnotesize]{$\tfrac12$}}
    edge from parent node[above, font=\footnotesize]{$\tfrac13$}
  };
\end{tikzpicture}
\caption{The Monty Hall probability tree (you have chosen Door $1$). Each path's probability
is the product of the edges along it. Only the two red leaves match the evidence ``Monty
opened Door $3$'': the $C_2$ world contributes weight $\tfrac13$, the $C_1$ world only
$\tfrac16$. Renormalizing these two ($\tfrac13 : \tfrac16 = 2 : 1$) gives posterior
$\P(C_2\given M_3)=\tfrac23$ for switching against $\P(C_1\given M_3)=\tfrac13$ for staying.
The asymmetry is born entirely in the forced ``$1$'' edge under $C_2$ versus the free
``$\tfrac12$'' edge under $C_1$.}
\label{ch05:fig-monty}
\end{figure}

\begin{happeningbox}
The crux is the asymmetric likelihood: $\P(M_3\given C_2)=1$ but $\P(M_3\given C_1)=\tfrac12$.
When the car is behind Door $2$, Monty is \emph{compelled} to open Door $3$; when the car
is behind your door, he had a free choice and opened Door $3$ only half the time. So the
observation ``Monty opened Door $3$'' is \emph{twice as likely} in the world where the car
is behind Door $2$. That factor of two is the Bayes factor, and it is exactly what shifts
your belief from $1/3$ to $2/3$. The information did not come from \emph{which} door Monty
opened so much as from the \emph{constraints he was under} when he chose --- constraints
created by his knowledge of the car's location.
\end{happeningbox}

\begin{trapbox}
Everything hinges on Monty's protocol, and changing it changes the answer. Suppose Monty
does \emph{not} know where the car is and opens one of the other doors \emph{at random},
and it happens to reveal a goat. Now $\P(M_3\given C_3)=\tfrac12$ instead of $0$ --- a
clueless Monty could have blundered into the car. Redo the arithmetic and you get
$\P(C_1\given M_3) = \P(C_2 \given M_3) = \tfrac12$: switching no longer helps! The
``$50/50$'' answer is correct --- for a \emph{different game}. The lesson is permanent:
\emph{a likelihood is a statement about the data-generating process}, and you cannot write
one down until you know the rules by which the evidence was produced. Most people get
Monty wrong because they never asked how Monty decides which door to open.
\end{trapbox}

\section{A First Taste of the Bayesian Worldview}
\label{ch05:worldview}

We have been using Bayes' theorem as a calculating device. But it is also the seed of an
entire philosophy of inference --- a way of doing statistics that differs from the
classical, ``frequentist'' approach you will meet later, and that has its own deep chapter
ahead (it returns in force when we study Bayesian inference). Let us plant the seed now.

The Bayesian stance is this: \emph{probability is a measure of belief}, and unknown
quantities --- not just future coin flips, but fixed-but-unknown facts like a drug's true
effectiveness or a physical constant --- can carry probability distributions that describe
your uncertainty about them. You start with a prior distribution over the unknown, you
observe data through a likelihood, and Bayes' theorem hands you a posterior distribution
that is your complete, updated state of knowledge. Inference is not a yes/no verdict; it
is a reshaping of a distribution.

\begin{aestheticbox}
What makes this worldview so attractive --- and so powerful in machine learning --- is
that it produces \emph{calibrated uncertainty}, not just point answers. A frequentist
estimate might say ``the click-through rate is $4.1\%$.'' A Bayesian posterior says ``here
is the full curve of how plausible each possible rate is, given what we've seen,'' and
from that single object you can read off a best guess, an honest interval, and the
probability of any statement you care about (``what's the chance the new design beats the
old by more than a point?''). When data are scarce, the prior gently regularizes your
conclusions; when data are abundant, the likelihood overwhelms the prior and Bayesians and
frequentists shake hands. The posterior is a complete accounting of what you know and how
sure you are --- which is exactly what a decision-maker needs.
\end{aestheticbox}

\begin{applicationbox}
\textbf{The naive Bayes spam filter --- building one by hand.}
For two decades the first line of defense in your inbox was a direct application of this
chapter. The hypotheses are $S = \text{spam}$ and $H = \text{ham}$ (legitimate mail); the
evidence is the set of words in the email. By the odds form \eqref{ch05:odds-form}, and
\emph{assuming words are conditionally independent given the class} (this is the ``naive''
assumption --- crude, since ``free'' and ``money'' surely co-occur, but astonishingly
effective),
\[
  \underbrace{\frac{\P(S\given \text{words})}{\P(H\given \text{words})}}_{\text{posterior odds}}
  = \underbrace{\frac{\P(S)}{\P(H)}}_{\text{prior odds}}
    \prod_{\text{word } w}\frac{\P(w\given S)}{\P(w\given H)}.
\]
Each word contributes a Bayes factor $\P(w\given S)/\P(w\given H)$, estimated by counting:
out of all words in known spam, what fraction were ``$w$''? Suppose training data give
$\P(\text{``free''}\given S)=0.05$ and $\P(\text{``free''}\given H)=0.001$. Then the word
``free'' carries a Bayes factor of $50$ --- seeing it multiplies the spam odds fifty-fold.
A word like ``meeting'' might have a factor of $1/20$, pulling toward ham. The filter
simply multiplies all the factors (in practice it \emph{adds the logs}, as in the
connection box, to avoid numerical underflow), starts from the base-rate prior odds of
spam, and flags the email if the posterior odds clear a threshold.

This is \emph{naive Bayes}, and beyond your inbox it classifies documents by topic,
diagnoses faults from sensor readings, and serves as the dependable baseline against which
fancier classifiers are measured. Two practical wrinkles complete it: \emph{Laplace
smoothing} (add a token count of $1$ to every word so a never-before-seen word doesn't
multiply your probability to a hard zero --- the dogmatic-prior trap again), and working
entirely in log-space. With those, a hundred lines of code and the theorem from page one
of this chapter gave the world a usable spam filter.
\end{applicationbox}

\begin{example}[A two-word spam score]
\label{ch05:ex-spam}
Take prior odds of spam $\P(S)/\P(H) = 0.4/0.6 = 2/3$ (your inbox is $40\%$ junk). An
email contains exactly the salient words ``free'' (Bayes factor $50$) and ``meeting''
(Bayes factor $1/20 = 0.05$). The posterior odds:
\[
  \frac{\P(S\given\text{words})}{\P(H\given\text{words})}
  = \frac23 \times 50 \times 0.05 = \frac23 \times 2.5 = \frac{5}{3} \approx 1.67.
\]
Posterior odds of $1.67{:}1$ means a posterior probability of
$1.67/(1+1.67)\approx 0.625$. The ``free'' shouted spam, the ``meeting'' whispered ham,
and on balance the message lands at about $63\%$ likely spam --- borderline, exactly the
kind of case a threshold has to adjudicate. Add one more spammy word and it tips over;
that is updating, word by word, the whole chapter in miniature.
\end{example}

\begin{labbox}
Before reading on, try these three to make sure the machinery is in your hands.
\begin{enumerate}
\item In Example~\ref{ch05:ex-rain}, suppose the streets are reported \emph{dry}. Without
peeking, is $\P(R \given \text{dry})$ bigger or smaller than the prior $0.2$? Compute it.
\item In the base-rate paradox, what disease prevalence $\P(D)$ would make a single
positive test a coin flip, i.e.\ $\P(D\given +)=0.5$? (Use the odds form: you need the
prior odds to equal $1/19.8$.)
\item In Monty Hall, redo the posterior if Monty, when he has a free choice (car behind
your door), prefers the \emph{lowest-numbered} available door. If he opens Door $3$, what
is $\P(C_1\given M_3)$ now? (This shows how a quirk in Monty's habits leaks extra
information.)
\end{enumerate}
\end{labbox}

\section{The Shape of What We Learned}
\label{ch05:summary}

Step back and notice how much reach came from how little. A single line of algebra ---
the symmetry of an intersection, divided through --- became a complete theory of how to
revise belief in the face of evidence.

\begin{itemize}
\item \textbf{Bayes' theorem} flips a conditional probability:
$\P(H\given E) = \P(E\given H)\P(H)/\P(E)$, turning the forward arrow ``cause $\to$
effect'' into the inferential arrow ``effect $\to$ likely cause.''
\item Its four parts are \textbf{prior} (belief before), \textbf{likelihood} (how the
hypothesis predicts the data), \textbf{evidence} (the normalizer, computed by the law of
total probability), and \textbf{posterior} (belief after). The slogan that carries it all:
\emph{posterior $\propto$ likelihood $\times$ prior}.
\item The \textbf{base-rate fallacy} is the failure to multiply by the prior: a reliable
detector of a rare thing fires mostly false alarms. Always ask how common the conclusion
was to begin with; count natural frequencies when the algebra feels slippery.
\item \textbf{Sequential updating} makes yesterday's posterior today's prior; with
conditionally independent evidence, beliefs compound by multiplying likelihoods --- or,
in the \textbf{odds form}, posterior odds $=$ Bayes factor $\times$ prior odds, which in
log-space is just addition.
\item \textbf{Monty Hall} is won by writing honest likelihoods --- which forces you to ask
how the evidence was generated. The host's knowledge is the whole secret.
\item The \textbf{Bayesian worldview} treats probability as belief and returns whole
posterior distributions, not just point estimates --- the foundation of \textbf{naive
Bayes} and a thread that runs all the way to modern machine learning.
\end{itemize}

What you should carry forward is less a formula than a habit of mind: when evidence
arrives, do not ask ``is this consistent with my hypothesis?'' --- almost everything is
consistent with almost everything. Ask instead, ``how much \emph{more} likely is this
evidence under my hypothesis than under its rivals, and how plausible was my hypothesis to
begin with?'' That question, asked honestly and repeatedly, is the entire art of reasoning
under uncertainty. The rest of this book gives it ever sharper tools.

\section{Exercises}
\label{ch05:exercises}

\begin{exercise}[Warming up the flip]
A factory has two machines. Machine A makes $60\%$ of the widgets and $2\%$ of its output
is defective; Machine B makes $40\%$ and $5\%$ is defective. A widget is drawn at random
and found defective. What is the probability it came from Machine B? (Find $\P(E)$ first,
then apply Bayes.)
\end{exercise}

\begin{exercise}[The base rate bites]
A new rapid test for a virus has sensitivity $0.98$ and specificity $0.96$. In a town
where $0.5\%$ of people are currently infected, a random resident tests positive. (a) Find
the posterior probability of infection. (b) The town doctor only tests people with
symptoms, among whom the prevalence is $30\%$. Recompute the posterior for a symptomatic
patient who tests positive, and explain in one sentence why the same test is now so much
more informative.
\end{exercise}

\begin{exercise}[Compounding evidence]
Using the odds form, show that for the test of Exercise~5.2(a) (prevalence $0.5\%$,
Bayes factor $0.98/0.04 = 24.5$), it takes how many \emph{independent} positive tests to
push the posterior probability of infection above $0.95$? (Hint: posterior odds $=
(1/199)\times 24.5^{\,n}$; solve for the smallest integer $n$ giving odds $\ge 19$.)
\end{exercise}

\begin{exercise}[Monty's quirk, in full]
Generalize the biased-host variant. When the car is behind your chosen Door $1$, Monty
opens Door $2$ with probability $p$ and Door $3$ with probability $1-p$ (when forced, he
opens whichever goat door he must). You see him open Door $3$. (a) Write the three
likelihoods $\P(M_3\given C_i)$ in terms of $p$. (b) Show that
$\P(C_2\given M_3) = 1/(2-p)$ and $\P(C_1\given M_3) = (1-p)/(2-p)$. (c) Verify that for
\emph{every} $p\in[0,1]$, switching is at least as good as staying, and find the $p$ that
makes staying \emph{least} bad.
\end{exercise}

\begin{exercise}[Three boxes, two coins]
Box 1 holds two gold coins, Box 2 one gold and one silver, Box 3 two silver. You pick a
box at random and draw one coin; it is gold. What is the probability the \emph{other}
coin in that box is also gold? (This is Bertrand's box paradox; the naive answer $1/2$ is
wrong --- find the right one and explain the error.)
\end{exercise}

\begin{exercise}[Build the filter]
You are training a tiny naive Bayes spam filter. Your corpus has $200$ spam and $300$ ham
emails. The word ``winner'' appears in $40$ of the spam and in $3$ of the ham.
(a) Estimate $\P(\text{``winner''}\given S)$ and $\P(\text{``winner''}\given H)$ and the
word's Bayes factor. (b) Using the prior odds from the corpus, find the posterior
probability of spam for an email whose only informative word is ``winner.'' (c) Explain
why Laplace smoothing (adding $1$ to each count) would matter for a word that appears in
spam but \emph{never} in your ham sample, and recompute the Bayes factor for such a word
that appears in $10$ spam and $0$ ham, with smoothing.
\end{exercise}

\begin{exercise}[The prosecutor's fallacy]
At a crime scene, DNA is recovered. A suspect's DNA matches, and an expert testifies that
``the probability of such a match in an innocent person is $1$ in a million.'' The
prosecutor concludes the suspect is guilty with probability $0.999999$. Suppose the
relevant population of possible perpetrators is $2$ million people, exactly one of whom is
guilty. (a) Using Bayes' theorem with a uniform prior over the $2$ million, find the
posterior probability of guilt given the match. (b) Explain in terms of this chapter
exactly which probabilities the prosecutor confused.
\end{exercise}

\begin{exercise}[Proof: order does not matter]
Let $H$ be a hypothesis and let $E_1, E_2$ be two pieces of evidence that are
conditionally independent given $H$ and given $H^c$. Prove, starting from the definition of
conditional probability, that
\[
  \P(H\given E_1, E_2) = \P(H \given E_2, E_1),
\]
and more strongly that the posterior depends on $(E_1,E_2)$ only through the product of
their individual Bayes factors. (This is the formal statement that conditionally
independent evidence can be processed in any order with the same result.)
\end{exercise}

\begin{exercise}[Conceptual: zero is forever]
Explain, with a short argument from Bayes' theorem, why a hypothesis assigned prior
probability exactly $0$ has posterior probability $0$ no matter what evidence arrives, and
why a hypothesis assigned prior $1$ stays at $1$. What practical advice for assigning
priors follows? (This is sometimes called Cromwell's rule.)
\end{exercise}

\begin{exercise}[A continuous preview]
A coin has unknown bias $\theta\in[0,1]$ of landing heads. Before flipping, you regard all
values of $\theta$ as equally likely (a uniform prior). You flip once and get heads.
Argue informally --- you may use the proportionality slogan --- that your posterior for
$\theta$ should now be proportional to $\theta$ (i.e.\ larger biases become more
plausible), and sketch its shape on $[0,1]$. (You have just done your first Bayesian
update on a \emph{continuous} parameter; the full treatment with the Beta distribution
awaits a later chapter.)
\end{exercise}
